% =====================================================================
%  THÈME 2 — Conversions en quantité de matière
%  Niveau DIFFICILE — Corrigé
% =====================================================================
\documentclass[11pt,a4paper]{article}
\usepackage[utf8]{inputenc}
\usepackage[T1]{fontenc}
\usepackage{lmodern}
\usepackage[french]{babel}
\usepackage{amsmath,amssymb,mathtools}
\usepackage{icomma}
\usepackage[version=4]{mhchem}
\usepackage{siunitx}
\sisetup{output-decimal-marker={,},per-mode=symbol}
\DeclareSIUnit\Molar{\mole\per\litre}
\usepackage{xcolor}
\usepackage{tikz}
\usepackage[most]{tcolorbox}
\usepackage{enumitem}
\usepackage{array}
\usepackage{fancyhdr}
\usepackage[a4paper,margin=2.1cm,headheight=26pt,headsep=8pt]{geometry}

\definecolor{primary}{RGB}{142,93,168}
\definecolor{primarydark}{RGB}{82,45,108}
\definecolor{corrcol}{RGB}{176,110,20}
\newcommand{\gmol}{\si{\gram\per\mole}}
\newcommand{\Vm}{V_{\mathrm m}}

\pagestyle{fancy}\fancyhf{}
\fancyhead[L]{\small\textcolor{primarydark}{\textbf{Fiche 5} — Thème 2 : Conversions}}
\fancyhead[R]{\small\textcolor{corrcol}{\textsc{Corrigé — Difficile}}}
\fancyfoot[C]{\small\thepage}
\renewcommand{\headrulewidth}{0.8pt}
\renewcommand{\headrule}{\hbox to\headwidth{\color{primary}\leaders\hrule height \headrulewidth\hfill}}

\newtcolorbox[auto counter]{cor}[1][]{enhanced,breakable,boxrule=1pt,arc=3pt,
  left=3mm,right=3mm,top=2.4mm,bottom=2.4mm,colback=corrcol!6,colframe=corrcol,
  fonttitle=\bfseries,coltitle=white,
  title={Corrigé~\thetcbcounter\ \ifx\relax#1\relax\relax\else\textmd{--- \itshape #1}\fi}}

\newcommand{\rep}[1]{\par\smallskip\noindent\textbf{\color{corrcol}$\Rightarrow$ #1}}

\setlength{\parindent}{0pt}
\setlength{\parskip}{3pt}

\begin{document}

\begin{center}
{\LARGE\bfseries\color{primarydark}Thème 2 — Conversions}\\[1pt]
{\large\color{corrcol}\textbf{Corrigé — Niveau Difficile}}\\
{\color{primary}\rule{0.6\linewidth}{1pt}}
\end{center}

\begin{cor}[identifier un métal alcalin par la masse molaire]
\textbf{a)} $M=\dfrac{m}{n}=\dfrac{42{,}4}{0{,}20}=\SI{212}{\gram\per\mole}$.\\
\textbf{b)} $M(\ce{X3PO4})=3\,M_{\ce{X}} + M(\ce{P}) + 4\,M(\ce{O}) = 3M_{\ce{X}} + 31 + 64 = 3M_{\ce{X}} + 95$.\\
\textbf{c)} $3M_{\ce{X}} + 95 = 212 \Rightarrow 3M_{\ce{X}} = 117 \Rightarrow M_{\ce{X}} = \SI{39}{\gram\per\mole}$.\\
\textbf{d)} $M_{\ce{X}}=39 \Rightarrow$ \boxed{\text{le potassium \ce{K}}} ; le composé est \ce{K3PO4}.
\rep{Méthode : masse molaire mesurée $\to$ équation en $M_{\ce{X}}$ $\to$ identification.}
\end{cor}

\begin{cor}[un sulfate métallique hydraté (mélantérite)]
\textbf{a)} $M=\dfrac{55{,}6}{0{,}20}=\SI{278}{\gram\per\mole}$.\\
\textbf{b)} $M(\ce{MSO4.7H2O})=M_{\ce{M}} + 32 + 4\times16 + 7\times18 = M_{\ce{M}} + 32 + 64 + 126 = M_{\ce{M}} + 222$.
Donc $M_{\ce{M}} = 278 - 222 = \SI{56}{\gram\per\mole}$.\\
\textbf{c)} $M_{\ce{M}}=56 \Rightarrow$ \boxed{\text{le fer \ce{Fe}}} ; le composé est \ce{FeSO4.7H2O}
(sulfate de fer(II) heptahydraté).
\end{cor}

\begin{cor}[du cube solide aux ions en solution]
\textbf{a)} Volume : $V=a^3=2^3=\SI{8}{\centi\meter\cubed}$.\quad Densité :
$\SI{2,0}{\kilo\gram\per\deci\meter\cubed} = \dfrac{\SI{2000}{\gram}}{\SI{1000}{\centi\meter\cubed}} = \SI{2,0}{\gram\per\centi\meter\cubed}$.\\
\textbf{b)} $m=\rho\,V = 2{,}0\times8 = \SI{16}{\gram}$.\\
\textbf{c)} $n(\ce{NaOH})=\dfrac{m}{M}=\dfrac{16}{40}=\SI{0,40}{\mole}$.\\
\textbf{d)} Chaque \ce{NaOH} libère un \ce{Na+} : $n(\ce{Na+})=\boxed{\SI{0,40}{\mole}}$.
\rep{Chaîne complète : volume $\to$ masse $\to$ mole $\to$ ions.}
\end{cor}

\begin{cor}[deux ions dans une eau minérale]
\textbf{a)} $\SI{90,0}{\milli\gram\per\litre}=\SI{0,0900}{\gram\per\litre}$ ;
$[\ce{Ca^2+}]=\dfrac{0{,}0900}{40}=\boxed{\SI{2,25e-3}{\Molar}}$.\\
\textbf{b)} $M(\ce{HCO3-})=1+12+3\times16=\SI{61}{\gram\per\mole}$ ;
$[\ce{HCO3-}]=\dfrac{0{,}122}{61}=\boxed{\SI{2,00e-3}{\Molar}}$.
\rep{Une concentration en masse (\si{\milli\gram\per\litre}) devient molaire en divisant par $M$.}
\end{cor}

\end{document}
